\documentclass[11pt,a4paper]{article}
% ============================================
% GIFT v3.4 — Shared Preamble
% ============================================
% Usage: % ============================================
% GIFT v3.4 — Shared Preamble
% ============================================
% Usage: % ============================================
% GIFT v3.4 — Shared Preamble
% ============================================
% Usage: \input{gift_preamble} after \documentclass

% ENCODING & FONTS
\usepackage[utf8]{inputenc}
\usepackage[T1]{fontenc}
\usepackage{lmodern}

% PAGE LAYOUT (Golden Ratio)
\usepackage[margin=1.618cm, top=2.618cm, bottom=2.618cm]{geometry}

% ESSENTIAL PACKAGES
\usepackage{float}
\usepackage{caption}
\usepackage{subcaption}
\usepackage{setspace}
\usepackage{fancyhdr}
\usepackage{xcolor}
\usepackage{hyperref}
\usepackage{csquotes}
\usepackage{amsmath}
\usepackage{amssymb}
\usepackage{amsthm}
\usepackage{mathtools}
\usepackage{booktabs}
\usepackage{longtable}
\usepackage{array}
\usepackage{multirow}
\usepackage{tikz}
\usepackage{graphicx}
\usepackage{listings}
\usepackage{enumitem}
% Unicode fallbacks (pandoc artifacts that survive markdown→LaTeX)
\DeclareUnicodeCharacter{00B0}{\ensuremath{^\circ}}
\DeclareUnicodeCharacter{221A}{\ensuremath{\sqrt{\,}}}
\DeclareUnicodeCharacter{2713}{\ensuremath{\checkmark}}
\DeclareUnicodeCharacter{2717}{\ensuremath{\times}}
\DeclareUnicodeCharacter{2070}{\ensuremath{^{0}}}
\DeclareUnicodeCharacter{00B9}{\ensuremath{^{1}}}
\DeclareUnicodeCharacter{00B2}{\ensuremath{^{2}}}
\DeclareUnicodeCharacter{00B3}{\ensuremath{^{3}}}
\DeclareUnicodeCharacter{2074}{\ensuremath{^{4}}}
\DeclareUnicodeCharacter{2075}{\ensuremath{^{5}}}
\DeclareUnicodeCharacter{2076}{\ensuremath{^{6}}}
\DeclareUnicodeCharacter{2077}{\ensuremath{^{7}}}
\DeclareUnicodeCharacter{2078}{\ensuremath{^{8}}}
\DeclareUnicodeCharacter{2079}{\ensuremath{^{9}}}
\DeclareUnicodeCharacter{2071}{\ensuremath{^{i}}}
\DeclareUnicodeCharacter{2080}{\ensuremath{_{0}}}
\DeclareUnicodeCharacter{2081}{\ensuremath{_{1}}}
\DeclareUnicodeCharacter{2082}{\ensuremath{_{2}}}
\DeclareUnicodeCharacter{2083}{\ensuremath{_{3}}}
\DeclareUnicodeCharacter{2084}{\ensuremath{_{4}}}
\DeclareUnicodeCharacter{2085}{\ensuremath{_{5}}}
\DeclareUnicodeCharacter{2086}{\ensuremath{_{6}}}
\DeclareUnicodeCharacter{2087}{\ensuremath{_{7}}}
\DeclareUnicodeCharacter{2088}{\ensuremath{_{8}}}
\DeclareUnicodeCharacter{2089}{\ensuremath{_{9}}}
\DeclareUnicodeCharacter{208A}{\ensuremath{_{+}}}
\DeclareUnicodeCharacter{208B}{\ensuremath{_{-}}}
\DeclareUnicodeCharacter{2192}{\ensuremath{\to}}
\DeclareUnicodeCharacter{2502}{|}
\DeclareUnicodeCharacter{25BC}{\ensuremath{\blacktriangledown}}
\DeclareUnicodeCharacter{2500}{\ensuremath{-}}
\DeclareUnicodeCharacter{2514}{\ensuremath{\llcorner}}
\DeclareUnicodeCharacter{251C}{\ensuremath{\vdash}}
\DeclareUnicodeCharacter{2026}{\ldots}
\DeclareUnicodeCharacter{2208}{\ensuremath{\in}}
\DeclareUnicodeCharacter{2229}{\ensuremath{\cap}}
\DeclareUnicodeCharacter{222A}{\ensuremath{\cup}}
\DeclareUnicodeCharacter{2295}{\ensuremath{\oplus}}
\DeclareUnicodeCharacter{2297}{\ensuremath{\otimes}}
\DeclareUnicodeCharacter{2245}{\ensuremath{\cong}}
\DeclareUnicodeCharacter{2248}{\ensuremath{\approx}}
\DeclareUnicodeCharacter{2260}{\ensuremath{\neq}}
\DeclareUnicodeCharacter{2261}{\ensuremath{\equiv}}
\DeclareUnicodeCharacter{2264}{\ensuremath{\leq}}
\DeclareUnicodeCharacter{2265}{\ensuremath{\geq}}
\DeclareUnicodeCharacter{00D7}{\ensuremath{\times}}
\DeclareUnicodeCharacter{00B1}{\ensuremath{\pm}}
\DeclareUnicodeCharacter{2225}{\ensuremath{\|}}
% Greek letters used in body text (pandoc passes them through unchanged)
\DeclareUnicodeCharacter{03B1}{\ensuremath{\alpha}}
\DeclareUnicodeCharacter{03B2}{\ensuremath{\beta}}
\DeclareUnicodeCharacter{03B3}{\ensuremath{\gamma}}
\DeclareUnicodeCharacter{03B4}{\ensuremath{\delta}}
\DeclareUnicodeCharacter{03B5}{\ensuremath{\varepsilon}}
\DeclareUnicodeCharacter{03B6}{\ensuremath{\zeta}}
\DeclareUnicodeCharacter{03B7}{\ensuremath{\eta}}
\DeclareUnicodeCharacter{03B8}{\ensuremath{\theta}}
\DeclareUnicodeCharacter{03B9}{\ensuremath{\iota}}
\DeclareUnicodeCharacter{03BA}{\ensuremath{\kappa}}
\DeclareUnicodeCharacter{03BB}{\ensuremath{\lambda}}
\DeclareUnicodeCharacter{03BC}{\ensuremath{\mu}}
\DeclareUnicodeCharacter{03BD}{\ensuremath{\nu}}
\DeclareUnicodeCharacter{03BE}{\ensuremath{\xi}}
\DeclareUnicodeCharacter{03C0}{\ensuremath{\pi}}
\DeclareUnicodeCharacter{03C1}{\ensuremath{\rho}}
\DeclareUnicodeCharacter{03C3}{\ensuremath{\sigma}}
\DeclareUnicodeCharacter{03C4}{\ensuremath{\tau}}
\DeclareUnicodeCharacter{03C5}{\ensuremath{\upsilon}}
\DeclareUnicodeCharacter{03C6}{\ensuremath{\varphi}}
\DeclareUnicodeCharacter{03C7}{\ensuremath{\chi}}
\DeclareUnicodeCharacter{03C8}{\ensuremath{\psi}}
\DeclareUnicodeCharacter{03C9}{\ensuremath{\omega}}
\DeclareUnicodeCharacter{0394}{\ensuremath{\Delta}}
\DeclareUnicodeCharacter{0398}{\ensuremath{\Theta}}
\DeclareUnicodeCharacter{039B}{\ensuremath{\Lambda}}
\DeclareUnicodeCharacter{03A0}{\ensuremath{\Pi}}
\DeclareUnicodeCharacter{03A3}{\ensuremath{\Sigma}}
\DeclareUnicodeCharacter{03A6}{\ensuremath{\Phi}}
\DeclareUnicodeCharacter{03A8}{\ensuremath{\Psi}}
\DeclareUnicodeCharacter{03A9}{\ensuremath{\Omega}}
\DeclareUnicodeCharacter{2200}{\ensuremath{\forall}}
\DeclareUnicodeCharacter{2203}{\ensuremath{\exists}}
\DeclareUnicodeCharacter{2207}{\ensuremath{\nabla}}
\DeclareUnicodeCharacter{2202}{\ensuremath{\partial}}
\DeclareUnicodeCharacter{2218}{\ensuremath{\circ}}
\DeclareUnicodeCharacter{2032}{\ensuremath{\prime}}
\DeclareUnicodeCharacter{2113}{\ensuremath{\ell}}
\DeclareUnicodeCharacter{210F}{\ensuremath{\hbar}}
\DeclareUnicodeCharacter{27E8}{\ensuremath{\langle}}
\DeclareUnicodeCharacter{27E9}{\ensuremath{\rangle}}
\DeclareUnicodeCharacter{2308}{\ensuremath{\lceil}}
\DeclareUnicodeCharacter{2309}{\ensuremath{\rceil}}
\DeclareUnicodeCharacter{230A}{\ensuremath{\lfloor}}
\DeclareUnicodeCharacter{230B}{\ensuremath{\rfloor}}
\DeclareUnicodeCharacter{222B}{\ensuremath{\int}}
\DeclareUnicodeCharacter{2211}{\ensuremath{\sum}}
\DeclareUnicodeCharacter{220F}{\ensuremath{\prod}}
\DeclareUnicodeCharacter{221E}{\ensuremath{\infty}}
\DeclareUnicodeCharacter{00B5}{\ensuremath{\mu}}
\DeclareUnicodeCharacter{2236}{\ensuremath{\colon}}
\DeclareUnicodeCharacter{2192}{\ensuremath{\to}}
\DeclareUnicodeCharacter{21A6}{\ensuremath{\mapsto}}
\DeclareUnicodeCharacter{21D2}{\ensuremath{\Rightarrow}}
\DeclareUnicodeCharacter{21D4}{\ensuremath{\Leftrightarrow}}
\DeclareUnicodeCharacter{2282}{\ensuremath{\subset}}
\DeclareUnicodeCharacter{2286}{\ensuremath{\subseteq}}
\DeclareUnicodeCharacter{2287}{\ensuremath{\supseteq}}
\DeclareUnicodeCharacter{222B}{\ensuremath{\int}}
\DeclareUnicodeCharacter{2205}{\ensuremath{\emptyset}}
\DeclareUnicodeCharacter{2299}{\ensuremath{\odot}}
\DeclareUnicodeCharacter{2220}{\ensuremath{\angle}}
\DeclareUnicodeCharacter{2135}{\ensuremath{\aleph}}
\DeclareUnicodeCharacter{207B}{\ensuremath{^{-}}}
\DeclareUnicodeCharacter{207A}{\ensuremath{^{+}}}
\DeclareUnicodeCharacter{220E}{\ensuremath{\blacksquare}}
\DeclareUnicodeCharacter{2061}{}
\DeclareUnicodeCharacter{2062}{}
\DeclareUnicodeCharacter{2063}{}
\DeclareUnicodeCharacter{2009}{\,}
\DeclareUnicodeCharacter{200A}{\,}
\DeclareUnicodeCharacter{200B}{}
\DeclareUnicodeCharacter{2212}{\ensuremath{-}}
\DeclareUnicodeCharacter{2194}{\ensuremath{\leftrightarrow}}
\DeclareUnicodeCharacter{2099}{\ensuremath{_{n}}}
\DeclareUnicodeCharacter{2098}{\ensuremath{_{m}}}
\DeclareUnicodeCharacter{2096}{\ensuremath{_{k}}}
\DeclareUnicodeCharacter{2095}{\ensuremath{_{h}}}
\DeclareUnicodeCharacter{2090}{\ensuremath{_{a}}}
\DeclareUnicodeCharacter{2091}{\ensuremath{_{e}}}
\DeclareUnicodeCharacter{2093}{\ensuremath{_{x}}}
\DeclareUnicodeCharacter{1D62}{\ensuremath{_{i}}}
\DeclareUnicodeCharacter{2C7C}{\ensuremath{_{j}}}
\DeclareUnicodeCharacter{2C7B}{\ensuremath{^{e}}}
\DeclareUnicodeCharacter{1D52}{\ensuremath{^{o}}}
\DeclareUnicodeCharacter{1D5}{\ensuremath{^{u}}}
\DeclareUnicodeCharacter{02E1}{\ensuremath{^{l}}}
\DeclareUnicodeCharacter{2295}{\ensuremath{\oplus}}
\DeclareUnicodeCharacter{27FA}{\ensuremath{\Longleftrightarrow}}
\DeclareUnicodeCharacter{27F9}{\ensuremath{\Longrightarrow}}
\DeclareUnicodeCharacter{2299}{\ensuremath{\odot}}
% Additional glyphs required by v3.5 markdown sources
\DeclareUnicodeCharacter{00A7}{\S}
\DeclareUnicodeCharacter{00B7}{\ensuremath{\cdot}}
\DeclareUnicodeCharacter{00E1}{\'{a}}
\DeclareUnicodeCharacter{00E4}{\"{a}}
\DeclareUnicodeCharacter{00E8}{\`{e}}
\DeclareUnicodeCharacter{00E9}{\'{e}}
\DeclareUnicodeCharacter{00F6}{\"{o}}
\DeclareUnicodeCharacter{010D}{\v{c}}
\DeclareUnicodeCharacter{014D}{\=o}
\DeclareUnicodeCharacter{0161}{\v{s}}
\DeclareUnicodeCharacter{0302}{}
\DeclareUnicodeCharacter{0304}{}
\DeclareUnicodeCharacter{2013}{--}
\DeclareUnicodeCharacter{2016}{\ensuremath{\|}}
\DeclareUnicodeCharacter{2020}{\ensuremath{\dagger}}
\DeclareUnicodeCharacter{2102}{\ensuremath{\mathbb{C}}}
\DeclareUnicodeCharacter{210D}{\ensuremath{\mathbb{H}}}
\DeclareUnicodeCharacter{211D}{\ensuremath{\mathbb{R}}}
\DeclareUnicodeCharacter{2124}{\ensuremath{\mathbb{Z}}}
\DeclareUnicodeCharacter{2190}{\ensuremath{\leftarrow}}
\DeclareUnicodeCharacter{2227}{\ensuremath{\wedge}}
\DeclareUnicodeCharacter{2228}{\ensuremath{\vee}}
\DeclareUnicodeCharacter{2283}{\ensuremath{\supset}}
\DeclareUnicodeCharacter{22C6}{\ensuremath{\star}}
\DeclareUnicodeCharacter{22EF}{\ensuremath{\cdots}}
\DeclareUnicodeCharacter{23F3}{\texttt{[\ldots]}}
% Provide \tightlist used by pandoc-generated lists (no-op)
\providecommand{\tightlist}{\setlength{\itemsep}{0pt}\setlength{\parskip}{0pt}}
% pandoc longtable column specs: provide \real as passthrough and a
% `none` counter so that calc's \real{x} expansion path works.
\newcounter{none}
\providecommand{\real}[1]{#1}

% LISTINGS
\lstset{
    basicstyle=\small\ttfamily,
    breaklines=true, frame=single,
    keepspaces=true, showstringspaces=false,
    aboveskip=0.8em, belowskip=0.8em
}

% HEADER/FOOTER
\setlength{\headheight}{14pt}
\pagestyle{fancy}
\fancyhf{}
\fancyhead[R]{\thepage}
\renewcommand{\headrulewidth}{0.2pt}

% HYPERREF
\hypersetup{
    colorlinks=true,
    linkcolor=blue, citecolor=blue, urlcolor=blue,
    pdfauthor={Brieuc de La Fourniere}
}

% SPACING
\setstretch{1.2}
\setlength{\parskip}{0.4em}
\setlength{\parindent}{0pt}

% TITLE FORMATTING
\usepackage{titling}
\pretitle{\LARGE\bfseries}
\posttitle{\vspace{-0.4em}}
\preauthor{}
\postauthor{}
\predate{}
\postdate{}
\setlength{\droptitle}{-2.0em}

% CUSTOM COMMANDS — GIFT notation
\newcommand{\E}{\mathrm{E}}
\newcommand{\Gtwo}{\mathrm{G}_2}
\newcommand{\Kseven}{K_7}
\newcommand{\AdS}{\mathrm{AdS}}
\newcommand{\SU}{\mathrm{SU}}
\newcommand{\SO}{\mathrm{SO}}
\newcommand{\U}{\mathrm{U}}
\newcommand{\Sp}{\mathrm{Sp}}
\newcommand{\PSL}{\mathrm{PSL}}
\newcommand{\SM}{\mathrm{SM}}
\newcommand{\Tr}{\mathrm{Tr}}
\newcommand{\rank}{\mathrm{rank}}
\newcommand{\Vol}{\mathrm{Vol}}
\newcommand{\Hol}{\mathrm{Hol}}
\newcommand{\Ric}{\mathrm{Ric}}
\newcommand{\Rm}{\mathrm{Rm}}
\newcommand{\Scal}{\mathrm{Scal}}
\DeclareMathOperator{\diag}{diag}
\DeclareMathOperator{\cond}{cond}
\newcommand{\GIFT}{\textrm{GIFT}}
\newcommand{\CP}{\mathrm{CP}}
\newcommand{\CKM}{\mathrm{CKM}}
\newcommand{\PMNS}{\mathrm{PMNS}}
\newcommand{\EW}{\mathrm{EW}}
\newcommand{\Pl}{\mathrm{Pl}}
\newcommand{\DE}{\mathrm{DE}}
\newcommand{\DM}{\mathrm{DM}}
\newcommand{\KK}{\mathrm{KK}}
\newcommand{\NK}{\mathrm{NK}}
\newcommand{\TCS}{\mathrm{TCS}}
\newcommand{\PINN}{\mathrm{PINN}}
\newcommand{\btwo}{b_2}
\newcommand{\bthree}{b_3}
\newcommand{\typeI}{\textsc{I}}
\newcommand{\typeII}{\textsc{II}}
\newcommand{\typeIII}{\textsc{III}}
\newcommand{\typeIV}{\textsc{IV}}
\newcommand{\proven}{\textsc{Proven}}

% PDF bookmark safety
\pdfstringdefDisableCommands{%
  \def\textsubscript#1{#1}%
  \def\textsuperscript#1{#1}%
  \def\CP{CP}%
  \def\Gtwo{G2}%
  \def\Kseven{K7}%
  \def\GIFT{GIFT}%
  \def\E{E}%
  \def\SM{SM}%
}

% THEOREM ENVIRONMENTS
\newtheorem{theorem}{Theorem}[section]
\newtheorem{proposition}[theorem]{Proposition}
\newtheorem{lemma}[theorem]{Lemma}
\newtheorem{corollary}[theorem]{Corollary}
\theoremstyle{definition}
\newtheorem{definition}[theorem]{Definition}
\theoremstyle{remark}
\newtheorem{remark}[theorem]{Remark}
 after \documentclass

% ENCODING & FONTS
\usepackage[utf8]{inputenc}
\usepackage[T1]{fontenc}
\usepackage{lmodern}

% PAGE LAYOUT (Golden Ratio)
\usepackage[margin=1.618cm, top=2.618cm, bottom=2.618cm]{geometry}

% ESSENTIAL PACKAGES
\usepackage{float}
\usepackage{caption}
\usepackage{subcaption}
\usepackage{setspace}
\usepackage{fancyhdr}
\usepackage{xcolor}
\usepackage{hyperref}
\usepackage{csquotes}
\usepackage{amsmath}
\usepackage{amssymb}
\usepackage{amsthm}
\usepackage{mathtools}
\usepackage{booktabs}
\usepackage{longtable}
\usepackage{array}
\usepackage{multirow}
\usepackage{tikz}
\usepackage{graphicx}
\usepackage{listings}
\usepackage{enumitem}
% Unicode fallbacks (pandoc artifacts that survive markdown→LaTeX)
\DeclareUnicodeCharacter{00B0}{\ensuremath{^\circ}}
\DeclareUnicodeCharacter{221A}{\ensuremath{\sqrt{\,}}}
\DeclareUnicodeCharacter{2713}{\ensuremath{\checkmark}}
\DeclareUnicodeCharacter{2717}{\ensuremath{\times}}
\DeclareUnicodeCharacter{2070}{\ensuremath{^{0}}}
\DeclareUnicodeCharacter{00B9}{\ensuremath{^{1}}}
\DeclareUnicodeCharacter{00B2}{\ensuremath{^{2}}}
\DeclareUnicodeCharacter{00B3}{\ensuremath{^{3}}}
\DeclareUnicodeCharacter{2074}{\ensuremath{^{4}}}
\DeclareUnicodeCharacter{2075}{\ensuremath{^{5}}}
\DeclareUnicodeCharacter{2076}{\ensuremath{^{6}}}
\DeclareUnicodeCharacter{2077}{\ensuremath{^{7}}}
\DeclareUnicodeCharacter{2078}{\ensuremath{^{8}}}
\DeclareUnicodeCharacter{2079}{\ensuremath{^{9}}}
\DeclareUnicodeCharacter{2071}{\ensuremath{^{i}}}
\DeclareUnicodeCharacter{2080}{\ensuremath{_{0}}}
\DeclareUnicodeCharacter{2081}{\ensuremath{_{1}}}
\DeclareUnicodeCharacter{2082}{\ensuremath{_{2}}}
\DeclareUnicodeCharacter{2083}{\ensuremath{_{3}}}
\DeclareUnicodeCharacter{2084}{\ensuremath{_{4}}}
\DeclareUnicodeCharacter{2085}{\ensuremath{_{5}}}
\DeclareUnicodeCharacter{2086}{\ensuremath{_{6}}}
\DeclareUnicodeCharacter{2087}{\ensuremath{_{7}}}
\DeclareUnicodeCharacter{2088}{\ensuremath{_{8}}}
\DeclareUnicodeCharacter{2089}{\ensuremath{_{9}}}
\DeclareUnicodeCharacter{208A}{\ensuremath{_{+}}}
\DeclareUnicodeCharacter{208B}{\ensuremath{_{-}}}
\DeclareUnicodeCharacter{2192}{\ensuremath{\to}}
\DeclareUnicodeCharacter{2502}{|}
\DeclareUnicodeCharacter{25BC}{\ensuremath{\blacktriangledown}}
\DeclareUnicodeCharacter{2500}{\ensuremath{-}}
\DeclareUnicodeCharacter{2514}{\ensuremath{\llcorner}}
\DeclareUnicodeCharacter{251C}{\ensuremath{\vdash}}
\DeclareUnicodeCharacter{2026}{\ldots}
\DeclareUnicodeCharacter{2208}{\ensuremath{\in}}
\DeclareUnicodeCharacter{2229}{\ensuremath{\cap}}
\DeclareUnicodeCharacter{222A}{\ensuremath{\cup}}
\DeclareUnicodeCharacter{2295}{\ensuremath{\oplus}}
\DeclareUnicodeCharacter{2297}{\ensuremath{\otimes}}
\DeclareUnicodeCharacter{2245}{\ensuremath{\cong}}
\DeclareUnicodeCharacter{2248}{\ensuremath{\approx}}
\DeclareUnicodeCharacter{2260}{\ensuremath{\neq}}
\DeclareUnicodeCharacter{2261}{\ensuremath{\equiv}}
\DeclareUnicodeCharacter{2264}{\ensuremath{\leq}}
\DeclareUnicodeCharacter{2265}{\ensuremath{\geq}}
\DeclareUnicodeCharacter{00D7}{\ensuremath{\times}}
\DeclareUnicodeCharacter{00B1}{\ensuremath{\pm}}
\DeclareUnicodeCharacter{2225}{\ensuremath{\|}}
% Greek letters used in body text (pandoc passes them through unchanged)
\DeclareUnicodeCharacter{03B1}{\ensuremath{\alpha}}
\DeclareUnicodeCharacter{03B2}{\ensuremath{\beta}}
\DeclareUnicodeCharacter{03B3}{\ensuremath{\gamma}}
\DeclareUnicodeCharacter{03B4}{\ensuremath{\delta}}
\DeclareUnicodeCharacter{03B5}{\ensuremath{\varepsilon}}
\DeclareUnicodeCharacter{03B6}{\ensuremath{\zeta}}
\DeclareUnicodeCharacter{03B7}{\ensuremath{\eta}}
\DeclareUnicodeCharacter{03B8}{\ensuremath{\theta}}
\DeclareUnicodeCharacter{03B9}{\ensuremath{\iota}}
\DeclareUnicodeCharacter{03BA}{\ensuremath{\kappa}}
\DeclareUnicodeCharacter{03BB}{\ensuremath{\lambda}}
\DeclareUnicodeCharacter{03BC}{\ensuremath{\mu}}
\DeclareUnicodeCharacter{03BD}{\ensuremath{\nu}}
\DeclareUnicodeCharacter{03BE}{\ensuremath{\xi}}
\DeclareUnicodeCharacter{03C0}{\ensuremath{\pi}}
\DeclareUnicodeCharacter{03C1}{\ensuremath{\rho}}
\DeclareUnicodeCharacter{03C3}{\ensuremath{\sigma}}
\DeclareUnicodeCharacter{03C4}{\ensuremath{\tau}}
\DeclareUnicodeCharacter{03C5}{\ensuremath{\upsilon}}
\DeclareUnicodeCharacter{03C6}{\ensuremath{\varphi}}
\DeclareUnicodeCharacter{03C7}{\ensuremath{\chi}}
\DeclareUnicodeCharacter{03C8}{\ensuremath{\psi}}
\DeclareUnicodeCharacter{03C9}{\ensuremath{\omega}}
\DeclareUnicodeCharacter{0394}{\ensuremath{\Delta}}
\DeclareUnicodeCharacter{0398}{\ensuremath{\Theta}}
\DeclareUnicodeCharacter{039B}{\ensuremath{\Lambda}}
\DeclareUnicodeCharacter{03A0}{\ensuremath{\Pi}}
\DeclareUnicodeCharacter{03A3}{\ensuremath{\Sigma}}
\DeclareUnicodeCharacter{03A6}{\ensuremath{\Phi}}
\DeclareUnicodeCharacter{03A8}{\ensuremath{\Psi}}
\DeclareUnicodeCharacter{03A9}{\ensuremath{\Omega}}
\DeclareUnicodeCharacter{2200}{\ensuremath{\forall}}
\DeclareUnicodeCharacter{2203}{\ensuremath{\exists}}
\DeclareUnicodeCharacter{2207}{\ensuremath{\nabla}}
\DeclareUnicodeCharacter{2202}{\ensuremath{\partial}}
\DeclareUnicodeCharacter{2218}{\ensuremath{\circ}}
\DeclareUnicodeCharacter{2032}{\ensuremath{\prime}}
\DeclareUnicodeCharacter{2113}{\ensuremath{\ell}}
\DeclareUnicodeCharacter{210F}{\ensuremath{\hbar}}
\DeclareUnicodeCharacter{27E8}{\ensuremath{\langle}}
\DeclareUnicodeCharacter{27E9}{\ensuremath{\rangle}}
\DeclareUnicodeCharacter{2308}{\ensuremath{\lceil}}
\DeclareUnicodeCharacter{2309}{\ensuremath{\rceil}}
\DeclareUnicodeCharacter{230A}{\ensuremath{\lfloor}}
\DeclareUnicodeCharacter{230B}{\ensuremath{\rfloor}}
\DeclareUnicodeCharacter{222B}{\ensuremath{\int}}
\DeclareUnicodeCharacter{2211}{\ensuremath{\sum}}
\DeclareUnicodeCharacter{220F}{\ensuremath{\prod}}
\DeclareUnicodeCharacter{221E}{\ensuremath{\infty}}
\DeclareUnicodeCharacter{00B5}{\ensuremath{\mu}}
\DeclareUnicodeCharacter{2236}{\ensuremath{\colon}}
\DeclareUnicodeCharacter{2192}{\ensuremath{\to}}
\DeclareUnicodeCharacter{21A6}{\ensuremath{\mapsto}}
\DeclareUnicodeCharacter{21D2}{\ensuremath{\Rightarrow}}
\DeclareUnicodeCharacter{21D4}{\ensuremath{\Leftrightarrow}}
\DeclareUnicodeCharacter{2282}{\ensuremath{\subset}}
\DeclareUnicodeCharacter{2286}{\ensuremath{\subseteq}}
\DeclareUnicodeCharacter{2287}{\ensuremath{\supseteq}}
\DeclareUnicodeCharacter{222B}{\ensuremath{\int}}
\DeclareUnicodeCharacter{2205}{\ensuremath{\emptyset}}
\DeclareUnicodeCharacter{2299}{\ensuremath{\odot}}
\DeclareUnicodeCharacter{2220}{\ensuremath{\angle}}
\DeclareUnicodeCharacter{2135}{\ensuremath{\aleph}}
\DeclareUnicodeCharacter{207B}{\ensuremath{^{-}}}
\DeclareUnicodeCharacter{207A}{\ensuremath{^{+}}}
\DeclareUnicodeCharacter{220E}{\ensuremath{\blacksquare}}
\DeclareUnicodeCharacter{2061}{}
\DeclareUnicodeCharacter{2062}{}
\DeclareUnicodeCharacter{2063}{}
\DeclareUnicodeCharacter{2009}{\,}
\DeclareUnicodeCharacter{200A}{\,}
\DeclareUnicodeCharacter{200B}{}
\DeclareUnicodeCharacter{2212}{\ensuremath{-}}
\DeclareUnicodeCharacter{2194}{\ensuremath{\leftrightarrow}}
\DeclareUnicodeCharacter{2099}{\ensuremath{_{n}}}
\DeclareUnicodeCharacter{2098}{\ensuremath{_{m}}}
\DeclareUnicodeCharacter{2096}{\ensuremath{_{k}}}
\DeclareUnicodeCharacter{2095}{\ensuremath{_{h}}}
\DeclareUnicodeCharacter{2090}{\ensuremath{_{a}}}
\DeclareUnicodeCharacter{2091}{\ensuremath{_{e}}}
\DeclareUnicodeCharacter{2093}{\ensuremath{_{x}}}
\DeclareUnicodeCharacter{1D62}{\ensuremath{_{i}}}
\DeclareUnicodeCharacter{2C7C}{\ensuremath{_{j}}}
\DeclareUnicodeCharacter{2C7B}{\ensuremath{^{e}}}
\DeclareUnicodeCharacter{1D52}{\ensuremath{^{o}}}
\DeclareUnicodeCharacter{1D5}{\ensuremath{^{u}}}
\DeclareUnicodeCharacter{02E1}{\ensuremath{^{l}}}
\DeclareUnicodeCharacter{2295}{\ensuremath{\oplus}}
\DeclareUnicodeCharacter{27FA}{\ensuremath{\Longleftrightarrow}}
\DeclareUnicodeCharacter{27F9}{\ensuremath{\Longrightarrow}}
\DeclareUnicodeCharacter{2299}{\ensuremath{\odot}}
% Additional glyphs required by v3.5 markdown sources
\DeclareUnicodeCharacter{00A7}{\S}
\DeclareUnicodeCharacter{00B7}{\ensuremath{\cdot}}
\DeclareUnicodeCharacter{00E1}{\'{a}}
\DeclareUnicodeCharacter{00E4}{\"{a}}
\DeclareUnicodeCharacter{00E8}{\`{e}}
\DeclareUnicodeCharacter{00E9}{\'{e}}
\DeclareUnicodeCharacter{00F6}{\"{o}}
\DeclareUnicodeCharacter{010D}{\v{c}}
\DeclareUnicodeCharacter{014D}{\=o}
\DeclareUnicodeCharacter{0161}{\v{s}}
\DeclareUnicodeCharacter{0302}{}
\DeclareUnicodeCharacter{0304}{}
\DeclareUnicodeCharacter{2013}{--}
\DeclareUnicodeCharacter{2016}{\ensuremath{\|}}
\DeclareUnicodeCharacter{2020}{\ensuremath{\dagger}}
\DeclareUnicodeCharacter{2102}{\ensuremath{\mathbb{C}}}
\DeclareUnicodeCharacter{210D}{\ensuremath{\mathbb{H}}}
\DeclareUnicodeCharacter{211D}{\ensuremath{\mathbb{R}}}
\DeclareUnicodeCharacter{2124}{\ensuremath{\mathbb{Z}}}
\DeclareUnicodeCharacter{2190}{\ensuremath{\leftarrow}}
\DeclareUnicodeCharacter{2227}{\ensuremath{\wedge}}
\DeclareUnicodeCharacter{2228}{\ensuremath{\vee}}
\DeclareUnicodeCharacter{2283}{\ensuremath{\supset}}
\DeclareUnicodeCharacter{22C6}{\ensuremath{\star}}
\DeclareUnicodeCharacter{22EF}{\ensuremath{\cdots}}
\DeclareUnicodeCharacter{23F3}{\texttt{[\ldots]}}
% Provide \tightlist used by pandoc-generated lists (no-op)
\providecommand{\tightlist}{\setlength{\itemsep}{0pt}\setlength{\parskip}{0pt}}
% pandoc longtable column specs: provide \real as passthrough and a
% `none` counter so that calc's \real{x} expansion path works.
\newcounter{none}
\providecommand{\real}[1]{#1}

% LISTINGS
\lstset{
    basicstyle=\small\ttfamily,
    breaklines=true, frame=single,
    keepspaces=true, showstringspaces=false,
    aboveskip=0.8em, belowskip=0.8em
}

% HEADER/FOOTER
\setlength{\headheight}{14pt}
\pagestyle{fancy}
\fancyhf{}
\fancyhead[R]{\thepage}
\renewcommand{\headrulewidth}{0.2pt}

% HYPERREF
\hypersetup{
    colorlinks=true,
    linkcolor=blue, citecolor=blue, urlcolor=blue,
    pdfauthor={Brieuc de La Fourniere}
}

% SPACING
\setstretch{1.2}
\setlength{\parskip}{0.4em}
\setlength{\parindent}{0pt}

% TITLE FORMATTING
\usepackage{titling}
\pretitle{\LARGE\bfseries}
\posttitle{\vspace{-0.4em}}
\preauthor{}
\postauthor{}
\predate{}
\postdate{}
\setlength{\droptitle}{-2.0em}

% CUSTOM COMMANDS — GIFT notation
\newcommand{\E}{\mathrm{E}}
\newcommand{\Gtwo}{\mathrm{G}_2}
\newcommand{\Kseven}{K_7}
\newcommand{\AdS}{\mathrm{AdS}}
\newcommand{\SU}{\mathrm{SU}}
\newcommand{\SO}{\mathrm{SO}}
\newcommand{\U}{\mathrm{U}}
\newcommand{\Sp}{\mathrm{Sp}}
\newcommand{\PSL}{\mathrm{PSL}}
\newcommand{\SM}{\mathrm{SM}}
\newcommand{\Tr}{\mathrm{Tr}}
\newcommand{\rank}{\mathrm{rank}}
\newcommand{\Vol}{\mathrm{Vol}}
\newcommand{\Hol}{\mathrm{Hol}}
\newcommand{\Ric}{\mathrm{Ric}}
\newcommand{\Rm}{\mathrm{Rm}}
\newcommand{\Scal}{\mathrm{Scal}}
\DeclareMathOperator{\diag}{diag}
\DeclareMathOperator{\cond}{cond}
\newcommand{\GIFT}{\textrm{GIFT}}
\newcommand{\CP}{\mathrm{CP}}
\newcommand{\CKM}{\mathrm{CKM}}
\newcommand{\PMNS}{\mathrm{PMNS}}
\newcommand{\EW}{\mathrm{EW}}
\newcommand{\Pl}{\mathrm{Pl}}
\newcommand{\DE}{\mathrm{DE}}
\newcommand{\DM}{\mathrm{DM}}
\newcommand{\KK}{\mathrm{KK}}
\newcommand{\NK}{\mathrm{NK}}
\newcommand{\TCS}{\mathrm{TCS}}
\newcommand{\PINN}{\mathrm{PINN}}
\newcommand{\btwo}{b_2}
\newcommand{\bthree}{b_3}
\newcommand{\typeI}{\textsc{I}}
\newcommand{\typeII}{\textsc{II}}
\newcommand{\typeIII}{\textsc{III}}
\newcommand{\typeIV}{\textsc{IV}}
\newcommand{\proven}{\textsc{Proven}}

% PDF bookmark safety
\pdfstringdefDisableCommands{%
  \def\textsubscript#1{#1}%
  \def\textsuperscript#1{#1}%
  \def\CP{CP}%
  \def\Gtwo{G2}%
  \def\Kseven{K7}%
  \def\GIFT{GIFT}%
  \def\E{E}%
  \def\SM{SM}%
}

% THEOREM ENVIRONMENTS
\newtheorem{theorem}{Theorem}[section]
\newtheorem{proposition}[theorem]{Proposition}
\newtheorem{lemma}[theorem]{Lemma}
\newtheorem{corollary}[theorem]{Corollary}
\theoremstyle{definition}
\newtheorem{definition}[theorem]{Definition}
\theoremstyle{remark}
\newtheorem{remark}[theorem]{Remark}
 after \documentclass

% ENCODING & FONTS
\usepackage[utf8]{inputenc}
\usepackage[T1]{fontenc}
\usepackage{lmodern}

% PAGE LAYOUT (Golden Ratio)
\usepackage[margin=1.618cm, top=2.618cm, bottom=2.618cm]{geometry}

% ESSENTIAL PACKAGES
\usepackage{float}
\usepackage{caption}
\usepackage{subcaption}
\usepackage{setspace}
\usepackage{fancyhdr}
\usepackage{xcolor}
\usepackage{hyperref}
\usepackage{csquotes}
\usepackage{amsmath}
\usepackage{amssymb}
\usepackage{amsthm}
\usepackage{mathtools}
\usepackage{booktabs}
\usepackage{longtable}
\usepackage{array}
\usepackage{multirow}
\usepackage{tikz}
\usepackage{graphicx}
\usepackage{listings}
\usepackage{enumitem}
% Unicode fallbacks (pandoc artifacts that survive markdown→LaTeX)
\DeclareUnicodeCharacter{00B0}{\ensuremath{^\circ}}
\DeclareUnicodeCharacter{221A}{\ensuremath{\sqrt{\,}}}
\DeclareUnicodeCharacter{2713}{\ensuremath{\checkmark}}
\DeclareUnicodeCharacter{2717}{\ensuremath{\times}}
\DeclareUnicodeCharacter{2070}{\ensuremath{^{0}}}
\DeclareUnicodeCharacter{00B9}{\ensuremath{^{1}}}
\DeclareUnicodeCharacter{00B2}{\ensuremath{^{2}}}
\DeclareUnicodeCharacter{00B3}{\ensuremath{^{3}}}
\DeclareUnicodeCharacter{2074}{\ensuremath{^{4}}}
\DeclareUnicodeCharacter{2075}{\ensuremath{^{5}}}
\DeclareUnicodeCharacter{2076}{\ensuremath{^{6}}}
\DeclareUnicodeCharacter{2077}{\ensuremath{^{7}}}
\DeclareUnicodeCharacter{2078}{\ensuremath{^{8}}}
\DeclareUnicodeCharacter{2079}{\ensuremath{^{9}}}
\DeclareUnicodeCharacter{2071}{\ensuremath{^{i}}}
\DeclareUnicodeCharacter{2080}{\ensuremath{_{0}}}
\DeclareUnicodeCharacter{2081}{\ensuremath{_{1}}}
\DeclareUnicodeCharacter{2082}{\ensuremath{_{2}}}
\DeclareUnicodeCharacter{2083}{\ensuremath{_{3}}}
\DeclareUnicodeCharacter{2084}{\ensuremath{_{4}}}
\DeclareUnicodeCharacter{2085}{\ensuremath{_{5}}}
\DeclareUnicodeCharacter{2086}{\ensuremath{_{6}}}
\DeclareUnicodeCharacter{2087}{\ensuremath{_{7}}}
\DeclareUnicodeCharacter{2088}{\ensuremath{_{8}}}
\DeclareUnicodeCharacter{2089}{\ensuremath{_{9}}}
\DeclareUnicodeCharacter{208A}{\ensuremath{_{+}}}
\DeclareUnicodeCharacter{208B}{\ensuremath{_{-}}}
\DeclareUnicodeCharacter{2192}{\ensuremath{\to}}
\DeclareUnicodeCharacter{2502}{|}
\DeclareUnicodeCharacter{25BC}{\ensuremath{\blacktriangledown}}
\DeclareUnicodeCharacter{2500}{\ensuremath{-}}
\DeclareUnicodeCharacter{2514}{\ensuremath{\llcorner}}
\DeclareUnicodeCharacter{251C}{\ensuremath{\vdash}}
\DeclareUnicodeCharacter{2026}{\ldots}
\DeclareUnicodeCharacter{2208}{\ensuremath{\in}}
\DeclareUnicodeCharacter{2229}{\ensuremath{\cap}}
\DeclareUnicodeCharacter{222A}{\ensuremath{\cup}}
\DeclareUnicodeCharacter{2295}{\ensuremath{\oplus}}
\DeclareUnicodeCharacter{2297}{\ensuremath{\otimes}}
\DeclareUnicodeCharacter{2245}{\ensuremath{\cong}}
\DeclareUnicodeCharacter{2248}{\ensuremath{\approx}}
\DeclareUnicodeCharacter{2260}{\ensuremath{\neq}}
\DeclareUnicodeCharacter{2261}{\ensuremath{\equiv}}
\DeclareUnicodeCharacter{2264}{\ensuremath{\leq}}
\DeclareUnicodeCharacter{2265}{\ensuremath{\geq}}
\DeclareUnicodeCharacter{00D7}{\ensuremath{\times}}
\DeclareUnicodeCharacter{00B1}{\ensuremath{\pm}}
\DeclareUnicodeCharacter{2225}{\ensuremath{\|}}
% Greek letters used in body text (pandoc passes them through unchanged)
\DeclareUnicodeCharacter{03B1}{\ensuremath{\alpha}}
\DeclareUnicodeCharacter{03B2}{\ensuremath{\beta}}
\DeclareUnicodeCharacter{03B3}{\ensuremath{\gamma}}
\DeclareUnicodeCharacter{03B4}{\ensuremath{\delta}}
\DeclareUnicodeCharacter{03B5}{\ensuremath{\varepsilon}}
\DeclareUnicodeCharacter{03B6}{\ensuremath{\zeta}}
\DeclareUnicodeCharacter{03B7}{\ensuremath{\eta}}
\DeclareUnicodeCharacter{03B8}{\ensuremath{\theta}}
\DeclareUnicodeCharacter{03B9}{\ensuremath{\iota}}
\DeclareUnicodeCharacter{03BA}{\ensuremath{\kappa}}
\DeclareUnicodeCharacter{03BB}{\ensuremath{\lambda}}
\DeclareUnicodeCharacter{03BC}{\ensuremath{\mu}}
\DeclareUnicodeCharacter{03BD}{\ensuremath{\nu}}
\DeclareUnicodeCharacter{03BE}{\ensuremath{\xi}}
\DeclareUnicodeCharacter{03C0}{\ensuremath{\pi}}
\DeclareUnicodeCharacter{03C1}{\ensuremath{\rho}}
\DeclareUnicodeCharacter{03C3}{\ensuremath{\sigma}}
\DeclareUnicodeCharacter{03C4}{\ensuremath{\tau}}
\DeclareUnicodeCharacter{03C5}{\ensuremath{\upsilon}}
\DeclareUnicodeCharacter{03C6}{\ensuremath{\varphi}}
\DeclareUnicodeCharacter{03C7}{\ensuremath{\chi}}
\DeclareUnicodeCharacter{03C8}{\ensuremath{\psi}}
\DeclareUnicodeCharacter{03C9}{\ensuremath{\omega}}
\DeclareUnicodeCharacter{0394}{\ensuremath{\Delta}}
\DeclareUnicodeCharacter{0398}{\ensuremath{\Theta}}
\DeclareUnicodeCharacter{039B}{\ensuremath{\Lambda}}
\DeclareUnicodeCharacter{03A0}{\ensuremath{\Pi}}
\DeclareUnicodeCharacter{03A3}{\ensuremath{\Sigma}}
\DeclareUnicodeCharacter{03A6}{\ensuremath{\Phi}}
\DeclareUnicodeCharacter{03A8}{\ensuremath{\Psi}}
\DeclareUnicodeCharacter{03A9}{\ensuremath{\Omega}}
\DeclareUnicodeCharacter{2200}{\ensuremath{\forall}}
\DeclareUnicodeCharacter{2203}{\ensuremath{\exists}}
\DeclareUnicodeCharacter{2207}{\ensuremath{\nabla}}
\DeclareUnicodeCharacter{2202}{\ensuremath{\partial}}
\DeclareUnicodeCharacter{2218}{\ensuremath{\circ}}
\DeclareUnicodeCharacter{2032}{\ensuremath{\prime}}
\DeclareUnicodeCharacter{2113}{\ensuremath{\ell}}
\DeclareUnicodeCharacter{210F}{\ensuremath{\hbar}}
\DeclareUnicodeCharacter{27E8}{\ensuremath{\langle}}
\DeclareUnicodeCharacter{27E9}{\ensuremath{\rangle}}
\DeclareUnicodeCharacter{2308}{\ensuremath{\lceil}}
\DeclareUnicodeCharacter{2309}{\ensuremath{\rceil}}
\DeclareUnicodeCharacter{230A}{\ensuremath{\lfloor}}
\DeclareUnicodeCharacter{230B}{\ensuremath{\rfloor}}
\DeclareUnicodeCharacter{222B}{\ensuremath{\int}}
\DeclareUnicodeCharacter{2211}{\ensuremath{\sum}}
\DeclareUnicodeCharacter{220F}{\ensuremath{\prod}}
\DeclareUnicodeCharacter{221E}{\ensuremath{\infty}}
\DeclareUnicodeCharacter{00B5}{\ensuremath{\mu}}
\DeclareUnicodeCharacter{2236}{\ensuremath{\colon}}
\DeclareUnicodeCharacter{2192}{\ensuremath{\to}}
\DeclareUnicodeCharacter{21A6}{\ensuremath{\mapsto}}
\DeclareUnicodeCharacter{21D2}{\ensuremath{\Rightarrow}}
\DeclareUnicodeCharacter{21D4}{\ensuremath{\Leftrightarrow}}
\DeclareUnicodeCharacter{2282}{\ensuremath{\subset}}
\DeclareUnicodeCharacter{2286}{\ensuremath{\subseteq}}
\DeclareUnicodeCharacter{2287}{\ensuremath{\supseteq}}
\DeclareUnicodeCharacter{222B}{\ensuremath{\int}}
\DeclareUnicodeCharacter{2205}{\ensuremath{\emptyset}}
\DeclareUnicodeCharacter{2299}{\ensuremath{\odot}}
\DeclareUnicodeCharacter{2220}{\ensuremath{\angle}}
\DeclareUnicodeCharacter{2135}{\ensuremath{\aleph}}
\DeclareUnicodeCharacter{207B}{\ensuremath{^{-}}}
\DeclareUnicodeCharacter{207A}{\ensuremath{^{+}}}
\DeclareUnicodeCharacter{220E}{\ensuremath{\blacksquare}}
\DeclareUnicodeCharacter{2061}{}
\DeclareUnicodeCharacter{2062}{}
\DeclareUnicodeCharacter{2063}{}
\DeclareUnicodeCharacter{2009}{\,}
\DeclareUnicodeCharacter{200A}{\,}
\DeclareUnicodeCharacter{200B}{}
\DeclareUnicodeCharacter{2212}{\ensuremath{-}}
\DeclareUnicodeCharacter{2194}{\ensuremath{\leftrightarrow}}
\DeclareUnicodeCharacter{2099}{\ensuremath{_{n}}}
\DeclareUnicodeCharacter{2098}{\ensuremath{_{m}}}
\DeclareUnicodeCharacter{2096}{\ensuremath{_{k}}}
\DeclareUnicodeCharacter{2095}{\ensuremath{_{h}}}
\DeclareUnicodeCharacter{2090}{\ensuremath{_{a}}}
\DeclareUnicodeCharacter{2091}{\ensuremath{_{e}}}
\DeclareUnicodeCharacter{2093}{\ensuremath{_{x}}}
\DeclareUnicodeCharacter{1D62}{\ensuremath{_{i}}}
\DeclareUnicodeCharacter{2C7C}{\ensuremath{_{j}}}
\DeclareUnicodeCharacter{2C7B}{\ensuremath{^{e}}}
\DeclareUnicodeCharacter{1D52}{\ensuremath{^{o}}}
\DeclareUnicodeCharacter{1D5}{\ensuremath{^{u}}}
\DeclareUnicodeCharacter{02E1}{\ensuremath{^{l}}}
\DeclareUnicodeCharacter{2295}{\ensuremath{\oplus}}
\DeclareUnicodeCharacter{27FA}{\ensuremath{\Longleftrightarrow}}
\DeclareUnicodeCharacter{27F9}{\ensuremath{\Longrightarrow}}
\DeclareUnicodeCharacter{2299}{\ensuremath{\odot}}
% Additional glyphs required by v3.5 markdown sources
\DeclareUnicodeCharacter{00A7}{\S}
\DeclareUnicodeCharacter{00B7}{\ensuremath{\cdot}}
\DeclareUnicodeCharacter{00E1}{\'{a}}
\DeclareUnicodeCharacter{00E4}{\"{a}}
\DeclareUnicodeCharacter{00E8}{\`{e}}
\DeclareUnicodeCharacter{00E9}{\'{e}}
\DeclareUnicodeCharacter{00F6}{\"{o}}
\DeclareUnicodeCharacter{010D}{\v{c}}
\DeclareUnicodeCharacter{014D}{\=o}
\DeclareUnicodeCharacter{0161}{\v{s}}
\DeclareUnicodeCharacter{0302}{}
\DeclareUnicodeCharacter{0304}{}
\DeclareUnicodeCharacter{2013}{--}
\DeclareUnicodeCharacter{2016}{\ensuremath{\|}}
\DeclareUnicodeCharacter{2020}{\ensuremath{\dagger}}
\DeclareUnicodeCharacter{2102}{\ensuremath{\mathbb{C}}}
\DeclareUnicodeCharacter{210D}{\ensuremath{\mathbb{H}}}
\DeclareUnicodeCharacter{211D}{\ensuremath{\mathbb{R}}}
\DeclareUnicodeCharacter{2124}{\ensuremath{\mathbb{Z}}}
\DeclareUnicodeCharacter{2190}{\ensuremath{\leftarrow}}
\DeclareUnicodeCharacter{2227}{\ensuremath{\wedge}}
\DeclareUnicodeCharacter{2228}{\ensuremath{\vee}}
\DeclareUnicodeCharacter{2283}{\ensuremath{\supset}}
\DeclareUnicodeCharacter{22C6}{\ensuremath{\star}}
\DeclareUnicodeCharacter{22EF}{\ensuremath{\cdots}}
\DeclareUnicodeCharacter{23F3}{\texttt{[\ldots]}}
% Provide \tightlist used by pandoc-generated lists (no-op)
\providecommand{\tightlist}{\setlength{\itemsep}{0pt}\setlength{\parskip}{0pt}}
% pandoc longtable column specs: provide \real as passthrough and a
% `none` counter so that calc's \real{x} expansion path works.
\newcounter{none}
\providecommand{\real}[1]{#1}

% LISTINGS
\lstset{
    basicstyle=\small\ttfamily,
    breaklines=true, frame=single,
    keepspaces=true, showstringspaces=false,
    aboveskip=0.8em, belowskip=0.8em
}

% HEADER/FOOTER
\setlength{\headheight}{14pt}
\pagestyle{fancy}
\fancyhf{}
\fancyhead[R]{\thepage}
\renewcommand{\headrulewidth}{0.2pt}

% HYPERREF
\hypersetup{
    colorlinks=true,
    linkcolor=blue, citecolor=blue, urlcolor=blue,
    pdfauthor={Brieuc de La Fourniere}
}

% SPACING
\setstretch{1.2}
\setlength{\parskip}{0.4em}
\setlength{\parindent}{0pt}

% TITLE FORMATTING
\usepackage{titling}
\pretitle{\LARGE\bfseries}
\posttitle{\vspace{-0.4em}}
\preauthor{}
\postauthor{}
\predate{}
\postdate{}
\setlength{\droptitle}{-2.0em}

% CUSTOM COMMANDS — GIFT notation
\newcommand{\E}{\mathrm{E}}
\newcommand{\Gtwo}{\mathrm{G}_2}
\newcommand{\Kseven}{K_7}
\newcommand{\AdS}{\mathrm{AdS}}
\newcommand{\SU}{\mathrm{SU}}
\newcommand{\SO}{\mathrm{SO}}
\newcommand{\U}{\mathrm{U}}
\newcommand{\Sp}{\mathrm{Sp}}
\newcommand{\PSL}{\mathrm{PSL}}
\newcommand{\SM}{\mathrm{SM}}
\newcommand{\Tr}{\mathrm{Tr}}
\newcommand{\rank}{\mathrm{rank}}
\newcommand{\Vol}{\mathrm{Vol}}
\newcommand{\Hol}{\mathrm{Hol}}
\newcommand{\Ric}{\mathrm{Ric}}
\newcommand{\Rm}{\mathrm{Rm}}
\newcommand{\Scal}{\mathrm{Scal}}
\DeclareMathOperator{\diag}{diag}
\DeclareMathOperator{\cond}{cond}
\newcommand{\GIFT}{\textrm{GIFT}}
\newcommand{\CP}{\mathrm{CP}}
\newcommand{\CKM}{\mathrm{CKM}}
\newcommand{\PMNS}{\mathrm{PMNS}}
\newcommand{\EW}{\mathrm{EW}}
\newcommand{\Pl}{\mathrm{Pl}}
\newcommand{\DE}{\mathrm{DE}}
\newcommand{\DM}{\mathrm{DM}}
\newcommand{\KK}{\mathrm{KK}}
\newcommand{\NK}{\mathrm{NK}}
\newcommand{\TCS}{\mathrm{TCS}}
\newcommand{\PINN}{\mathrm{PINN}}
\newcommand{\btwo}{b_2}
\newcommand{\bthree}{b_3}
\newcommand{\typeI}{\textsc{I}}
\newcommand{\typeII}{\textsc{II}}
\newcommand{\typeIII}{\textsc{III}}
\newcommand{\typeIV}{\textsc{IV}}
\newcommand{\proven}{\textsc{Proven}}

% PDF bookmark safety
\pdfstringdefDisableCommands{%
  \def\textsubscript#1{#1}%
  \def\textsuperscript#1{#1}%
  \def\CP{CP}%
  \def\Gtwo{G2}%
  \def\Kseven{K7}%
  \def\GIFT{GIFT}%
  \def\E{E}%
  \def\SM{SM}%
}

% THEOREM ENVIRONMENTS
\newtheorem{theorem}{Theorem}[section]
\newtheorem{proposition}[theorem]{Proposition}
\newtheorem{lemma}[theorem]{Lemma}
\newtheorem{corollary}[theorem]{Corollary}
\theoremstyle{definition}
\newtheorem{definition}[theorem]{Definition}
\theoremstyle{remark}
\newtheorem{remark}[theorem]{Remark}

% Glyphs new in 3.6, kept here rather than in the preamble shared with 3.4 and 3.5.
\DeclareUnicodeCharacter{02E2}{\ensuremath{^{\mathrm{S}}}}
\DeclareUnicodeCharacter{1D37}{\ensuremath{^{\mathrm{K}}}}
\DeclareUnicodeCharacter{1D39}{\ensuremath{^{\mathrm{M}}}}
\DeclareUnicodeCharacter{1D3A}{\ensuremath{^{\mathrm{N}}}}
\DeclareUnicodeCharacter{1D3E}{\ensuremath{^{\mathrm{P}}}}
\DeclareUnicodeCharacter{1D9C}{\ensuremath{^{\mathrm{C}}}}
\DeclareUnicodeCharacter{211A}{\ensuremath{\mathbb{Q}}}
\DeclareUnicodeCharacter{226B}{\ensuremath{\gg}}
\DeclareUnicodeCharacter{22A5}{\ensuremath{\perp}}
% pandoc >= 3.2.1 emits \pandocbounded around images; provide a fallback.
% Also clamp all images to \linewidth (natural-size figures overflowed the
% right margin -- the "truncated figures" flagged by the round-2 reviews).
% Kept here rather than in gift_preamble.tex, shared with the 3.4 build.
\providecommand{\pandocbounded}[1]{#1}
\makeatletter
\def\maxwidth{\ifdim\Gin@nat@width>\linewidth\linewidth\else\Gin@nat@width\fi}
\makeatother
\setkeys{Gin}{width=\maxwidth,keepaspectratio}
% md headings carry their own literal numbering ("3.2 What ..."); suppress
% the LaTeX section counters to avoid "0.2 1. Introduction" double numbering.
\setcounter{secnumdepth}{0}
% Loosen justification slightly: long technical identifiers (made breakable
% at build time via \allowbreak) still need extra stretch to avoid overfulls.
\emergencystretch=3em
\fancyhead[L]{The $K_7$ Dictionary v3.6: Supplement S4}
\hypersetup{pdftitle={The K7 Dictionary v3.6 Supplement S4: Sieve and Look-Elsewhere Diagnostics}}

\graphicspath{{../../../figures/}}

\title{%
\LARGE\textbf{Supplement S4: Sieve and Look-Elsewhere Diagnostics}\\[0.3em]
\large The Look-Elsewhere Null and Descriptive Diagnostics%
}
\author{Brieuc de La Fourni\`ere \\ \textit{Independent researcher, Beaune, France}}
\date{November 2026 \quad\textbar\quad v3.6}

\begin{document}

\maketitle
\noindent\rule{\textwidth}{0.2pt}

\thispagestyle{fancy}

\emph{Companion to the K₇ framework paper v3.6. This supplement gives the details of the two look-elsewhere readings of main §7: the published Sieve scorecard (S4.1) and the preregistered null in which every competitor searches for its own formulas (S4.2). It then archives, with their status, the descriptive diagnostics of version 3.5 (S4.3) and the coincidence probabilities of version 3.4 (S4.4). Neither archive is evidence for the framework. The tables of S4.2 are generated from the result files of the test and are not edited by hand.}

\section{S4.1 The Sieve scorecard}\label{s4.1-the-sieve-scorecard}

The Sieve {[}M{]} is a public battery of four nulls with a frozen observable list and grammar {[}S-frozen{]}. N1 asks whether a relation is isolated in the fixed vocabulary; N2 prices the freedom to have chosen the vocabulary, by comparing with budget-matched invented vocabularies; N3 and N4 shuffle coefficients and targets. Its published scorecard for the K₇ relations {[}M, §4.3{]} reads:

\begin{itemize}
\tightlist
\item
  n\textsubscript{s} and m\textsubscript{H}/m\textsubscript{W} survive N1 on all three of its axes.
\item
  Neither clears N2: a budget-matched invented vocabulary matches them as well (for m\textsubscript{H}/m\textsubscript{W}, 22 of 24 invented vocabularies do).
\item
  Q\textsubscript{Koide} is isolated in the fixed vocabulary, but its value 2/3 is reached by every vocabulary: it is a tie, not a survivor.
\item
  Across the nine shallow relations tested by N2, the invented-vocabulary envelope is beaten 0 times out of 9.
\end{itemize}

Version 3.5 of this supplement (§S4.4 there) described m\textsubscript{H}/m\textsubscript{W} as the ``unique survivor at exact rank 1'', and Koide and n\textsubscript{s} as passing the battery. That reading left out the N2 result of the same scorecard and is withdrawn. What N2 cannot price is where the vocabulary comes from: the K₇ constants are E₈ dimensions and ranks and the Betti numbers, not an invented list. That is a fact about the origin of the constants, not evidence that they were selected before the data were known: the grammar and the observable list were frozen after the formulas had been found (main §3.2), and main §10 has not justified the Betti pair at the selecting step.

\section{S4.2 A null in which every competitor searches for its own formulas}\label{s4.2-a-null-in-which-every-competitor-searches-for-its-own-formulas}

\subsection{S4.2.1 Question and design}\label{s4.2.1-question-and-design}

Version 3.5 stated that among 192 349 alternative configurations, none outperformed the framework (P \textless{} 5 × 10⁻⁶). The Monte Carlo behind that figure kept the 33 formulas of K₇ fixed and re-evaluated them with the invariants of each alternative; no alternative searched for its own formulas. Since the formulas had been chosen for (21, 77), (21, 77) won by construction. Several of the 33 predictions did not depend on the configuration at all, and 1/192 349 was simply the floor of the estimator. The figure is withdrawn.

The test that replaces it asks: \textbf{does K₇ remain exceptional when each competitor searches for its own formulas, in the same grammar, at the same budget, with the same freedom of assignment and of unit?} It was preregistered on 2026-10-17, before any run, with two outcomes only: a calibrated global P-value, or withdrawal of the P-value from the text.

\begin{itemize}
\tightlist
\item
  \textbf{Targets.} The 33 Type I observables of the dataset (S3), with their measured values and uncertainties.
\item
  \textbf{Grammar.} The frozen Sieve engine {[}M{]}: the leaves of its transcendental grammar (1 to 9, π, e, φ, γ, ln 2, ζ(3), ζ(5), ζ(7), ζ(9), ζ(11)) plus the 19 integer constants of the competitor's register, so that every competitor has 38 leaves (a constant alphabet size, which keeps a larger vocabulary from winning by size alone); the operations + − × / \^{}, square root, inverse, square and logarithm, with the engine's guards; \textbf{at most five nodes} per expression.
\item
  \textbf{Concordance.} A target is reached by a vocabulary if one of its expressions falls in the target's window. The window is {[}x − 2σ₋, x + 2σ₊{]} around the measured value x, with asymmetric uncertainties kept asymmetric (the Sieve's frozen rule τ\textsubscript{z} = 2); for m\textsubscript{s}/m\textsubscript{d}, which the PDG gives as a range with no uncertainty, the window is that range. For δ\textsubscript{CP}, the target is reached if the phase is reached in any of four units: degrees, radians, grads or turns.
\item
  \textbf{Assignment.} The search is blind: each target takes the cheapest concordant expression of the vocabulary. The full freedom to assign formulas to observables is therefore granted to everyone, \textbf{K₇ included}: K₇ is scored by the same blind search on its own vocabulary, not by its published formulas.
\item
  \textbf{Scores.} S\textsubscript{cover}, the number of targets reached (primary). S\textsubscript{cost}, the sum over targets of the description length of the cheapest concordant expression, an unreached target costing the largest attainable description length plus one operation (dual; reported, outside the rule).
\item
  \textbf{K₇'s register as tested.} The nineteen integers 2, 3, 5, 7, 8, 11, 14, 21, 27, 32, 52, 56, 61, 65, 77, 78, 99, 168, 248, plus φ: the register of the version 3.5 Monte Carlo (\texttt{algebraic\_\allowbreak{}montecarlo.\allowbreak{}py}), not literally the twenty constants of main §3.1. Relative to §3.1, dim E₇ = 133 and α\textsubscript{sum} = 13 are absent, and κ\textsubscript{T}⁻¹ = 61 is present. At a fixed budget the difference matters (13 costs two leaves, 8 + 5); a test on the literal register of §3.1 is a separate computation, not done here.
\item
  \textbf{Competitors.} F-BETTI: every pair b₂ ∈ {[}1, 40{]}, b₃ ∈ {[}17, 120{]} other than (21, 77), 4 159 pairs. Only H*, κ\textsubscript{T}⁻¹, det\textsubscript{den} and det\textsubscript{num} are re-derived, by the published definitions (main App. B): H* = b₂ + b₃ + 1, κ\textsubscript{T}⁻¹ = b₃ − dim(G₂) − p₂, det(g) = p₂ + 1/(b₂ + dim(G₂) − N\textsubscript{gen}). Every other constant is held at K₇'s value, N\textsubscript{gen} = 3 included, although main §3.1 writes N\textsubscript{gen} = b₂/dim K₇: the family is conditional on that choice. F-VOCAB: 2 000 vocabularies of 19 integers drawn by the Sieve's frozen generator (without replacement in 3 to 248, seed 20261017), plus φ.
\item
  \textbf{Verdict.} For each family F of size N, P\textsubscript{F} = (1 + n)/(1 + N), where n is the number of competitors with S\textsubscript{cover} at least that of K₇ (ties count against K₇). P\textsubscript{global} = max(P\textsubscript{F-BETTI}, P\textsubscript{F-VOCAB}): K₇ must beat both reference classes. The look-elsewhere effect is overcome if P\textsubscript{global} ≤ 0.01; otherwise the P-value is withdrawn (LEE-RETRAIT). Neither S\textsubscript{cost} nor the variant below can change the verdict.
\item
  \textbf{Gates.} G1 replays the engine on the (21, 77) row of the earlier landscape scan, target by target, to 10⁻⁹. G2 computes for each target whether K₇'s published value is in the set of its own expressions of at most five nodes, and requires that every such value within its window be reached by the blind search. If either fails, nothing is read.
\item
  \textbf{Variant, outside the rule.} Two derivations of κ\textsubscript{T}⁻¹ and det(g) circulate in the project: the published one above, and an earlier one (κ\textsubscript{T}⁻¹ = b₃ − b₂ + Weyl, det(g) = Weyl·(N\textsubscript{gen} + rank(E₈) + p₂)/2\textsuperscript{Weyl}) that agrees with it on the line b₂ = 21 and differs elsewhere. Which one is used elsewhere is itself a freedom of the dictionary; the variant family (4 160 pairs, (21, 77) included) reports its effect.
\end{itemize}

The preregistration was amended twice. The first amendment (2026-10-22) came after the external review of the v3.6 draft: the dataset was refreshed (S3) with all uncertainties brought to 1σ, the windows were made two-sided (the first run used only the side of each measurement that faces K₇'s value, which biased the search of the competitors towards that value), and the PDG range became the window of m\textsubscript{s}/m\textsubscript{d}.~The second (2026-11-03), after a further external review and before the current run, replaced the windows of the six Planck ratios by those computed on the official chains (S3; the window of Ω\textsubscript{c}/Ω\textsubscript{m} narrows by a factor of six), and added the witness below. In neither case were the question, families, grammar, budget, seed, τ\textsubscript{z} or verdict rule changed.

\subsection{S4.2.2 Result}\label{s4.2.2-result}

Both gates pass: G1 reproduces the 28 landscape-scan rows with no disagreement, and G2 finds 28 of K₇'s 33 published values among its own expressions of at most five nodes, every one of them within its window being reached. G3 count: 24 of the 33 published formulas lie within their windows (the 23 within 2σ of main §4, and m\textsubscript{s}/m\textsubscript{d} inside its PDG range); this is a diagnostic, not a statistic, since these formulas were chosen for (21, 77).

{\footnotesize\setlength{\tabcolsep}{3pt}
\begin{longtable}[]{@{}
  >{\raggedright\arraybackslash}p{(\linewidth - 12\tabcolsep) * \real{0.1111}}
  >{\raggedleft\arraybackslash}p{(\linewidth - 12\tabcolsep) * \real{0.1481}}
  >{\raggedleft\arraybackslash}p{(\linewidth - 12\tabcolsep) * \real{0.1481}}
  >{\raggedleft\arraybackslash}p{(\linewidth - 12\tabcolsep) * \real{0.1481}}
  >{\raggedleft\arraybackslash}p{(\linewidth - 12\tabcolsep) * \real{0.1481}}
  >{\raggedleft\arraybackslash}p{(\linewidth - 12\tabcolsep) * \real{0.1481}}
  >{\raggedleft\arraybackslash}p{(\linewidth - 12\tabcolsep) * \real{0.1481}}@{}}
\toprule\noalign{}
\begin{minipage}[b]{\linewidth}\raggedright
Family
\end{minipage} & \begin{minipage}[b]{\linewidth}\raggedleft
N
\end{minipage} & \begin{minipage}[b]{\linewidth}\raggedleft
reach ≥ 31/33 (as many as K₇)
\end{minipage} & \begin{minipage}[b]{\linewidth}\raggedleft
P = (1 + n)/(1 + N)
\end{minipage} & \begin{minipage}[b]{\linewidth}\raggedleft
max S\textsubscript{cover}
\end{minipage} & \begin{minipage}[b]{\linewidth}\raggedleft
S\textsubscript{cost} range (bits)
\end{minipage} & \begin{minipage}[b]{\linewidth}\raggedleft
cost ≤ K₇
\end{minipage} \\
\midrule\noalign{}
\endhead
\bottomrule\noalign{}
\endlastfoot
K₇ (21, 77), blind search on its own vocabulary & 1 & 31/33 & n/a & 31 & 408.1 & n/a \\
F-BETTI: alternative Betti pairs (published derivations) & 4 159 & 1 085 & 0.261 & 31 & 389.2 to 418.7 & 1 163 \\
F-VOCAB: invented vocabularies (Sieve generator) & 2 000 & 1 015 & 0.508 & 31 & 385.0 to 430.4 & 698 \\
F-BETTI-FICHE: alternative Betti pairs, variant derivations (outside the rule) & 4 160 & 723 & 0.174 & 31 & 391.3 to 417.6 & 601 \\
\end{longtable}}

P\textsubscript{global} = max(P\textsubscript{F-BETTI}, P\textsubscript{F-VOCAB}) = \textbf{0.508}; threshold 0.01; verdict \textbf{LEE-RETRAIT}. The variant row does not enter the verdict. ``Cost ≤ K₇'' counts floating-point ties as ties (\textbar ΔS\textsubscript{cost}\textbar{} ≤ 1e-09 bit; every tie is at the same (length, operations) pair as K₇).

Per target, the blind search of K₇ and the frequency with which each family reaches the target:

{\footnotesize\setlength{\tabcolsep}{3pt}
\begin{longtable}[]{@{}
  >{\raggedright\arraybackslash}p{(\linewidth - 12\tabcolsep) * \real{0.1304}}
  >{\raggedright\arraybackslash}p{(\linewidth - 12\tabcolsep) * \real{0.1304}}
  >{\raggedright\arraybackslash}p{(\linewidth - 12\tabcolsep) * \real{0.1304}}
  >{\raggedright\arraybackslash}p{(\linewidth - 12\tabcolsep) * \real{0.1304}}
  >{\raggedright\arraybackslash}p{(\linewidth - 12\tabcolsep) * \real{0.1304}}
  >{\raggedleft\arraybackslash}p{(\linewidth - 12\tabcolsep) * \real{0.1739}}
  >{\raggedleft\arraybackslash}p{(\linewidth - 12\tabcolsep) * \real{0.1739}}@{}}
\toprule\noalign{}
\begin{minipage}[b]{\linewidth}\raggedright
ID
\end{minipage} & \begin{minipage}[b]{\linewidth}\raggedright
Observable
\end{minipage} & \begin{minipage}[b]{\linewidth}\raggedright
K₇ reaches (blind)
\end{minipage} & \begin{minipage}[b]{\linewidth}\raggedright
published value in K₇'s ≤ 5-node set
\end{minipage} & \begin{minipage}[b]{\linewidth}\raggedright
published value within window
\end{minipage} & \begin{minipage}[b]{\linewidth}\raggedleft
Betti pairs reaching (of 4 159)
\end{minipage} & \begin{minipage}[b]{\linewidth}\raggedleft
vocabularies reaching (of 2 000)
\end{minipage} \\
\midrule\noalign{}
\endhead
\bottomrule\noalign{}
\endlastfoot
A01 & sin²θ\textsubscript{W} & yes & yes & no & 4 159 & 2 000 \\
A02 & α\textsubscript{s}(M\textsubscript{Z}) & yes & yes & yes & 4 159 & 2 000 \\
A03 & λ\textsubscript{H} & yes & no & yes & 4 159 & 2 000 \\
A04 & α⁻¹(0) & no & no & no & 0 & 0 \\
A05 & Q\textsubscript{Koide} & yes & yes & yes & 4 159 & 2 000 \\
A06 & m\textsubscript{τ}/m\textsubscript{e} & yes & yes & no & 1 085 & 1 015 \\
A07 & m\textsubscript{μ}/m\textsubscript{e} & no & no & no & 0 & 0 \\
A08 & m\textsubscript{s}/m\textsubscript{d} & yes & yes & yes & 4 159 & 2 000 \\
A09 & m\textsubscript{c}/m\textsubscript{s} & yes & yes & yes & 4 159 & 2 000 \\
A10 & m\textsubscript{b}/m\textsubscript{t} & yes & yes & no & 4 159 & 2 000 \\
A11 & m\textsubscript{u}/m\textsubscript{d} & yes & yes & yes & 4 159 & 2 000 \\
A12 & sin²θ₁₂ᴾᴹᴺˢ & yes & yes & yes & 4 159 & 2 000 \\
A13 & sin²θ₂₃ᴾᴹᴺˢ & yes & yes & no & 4 159 & 2 000 \\
A14 & sin²θ₁₃ᴾᴹᴺˢ & yes & yes & yes & 4 159 & 2 000 \\
A15 & δ\textsubscript{CP} & yes & yes & yes & 4 159 & 2 000 \\
A16 & sin θ₁₂ᶜᴷᴹ & yes & yes & yes & 4 159 & 2 000 \\
A17 & A\textsubscript{Wolf} & yes & yes & yes & 4 159 & 2 000 \\
A18 & sin θ₂₃ᶜᴷᴹ & yes & yes & yes & 4 159 & 2 000 \\
A19 & m\textsubscript{H}/m\textsubscript{t} & yes & yes & yes & 4 159 & 2 000 \\
A20 & m\textsubscript{H}/m\textsubscript{W} & yes & yes & yes & 4 159 & 2 000 \\
A21 & m\textsubscript{W}/m\textsubscript{Z} & yes & no & no & 4 159 & 2 000 \\
A22 & m\textsubscript{μ}/m\textsubscript{τ} & yes & yes & no & 4 159 & 2 000 \\
A23 & Ω\textsubscript{DE} & yes & yes & yes & 4 159 & 2 000 \\
A24 & n\textsubscript{s} & yes & no & yes & 4 159 & 2 000 \\
A25 & Ω\textsubscript{DM}/Ω\textsubscript{b} & yes & yes & yes & 4 159 & 2 000 \\
A26 & Ω\textsubscript{c}/Ω\textsubscript{Λ} & yes & yes & yes & 4 159 & 2 000 \\
A27 & Ω\textsubscript{Λ}/Ω\textsubscript{m} & yes & yes & yes & 4 159 & 2 000 \\
A28 & h & yes & yes & yes & 4 159 & 2 000 \\
A29 & Ω\textsubscript{b}/Ω\textsubscript{m} & yes & yes & yes & 4 159 & 2 000 \\
A30 & Ω\textsubscript{c}/Ω\textsubscript{m} & yes & yes & no & 4 159 & 2 000 \\
A31 & σ\textsubscript{8} & yes & yes & yes & 4 159 & 2 000 \\
A32 & Y\textsubscript{p} & yes & yes & yes & 4 159 & 2 000 \\
A33 & Ω\textsubscript{m}/Ω\textsubscript{Λ} & yes & yes & yes & 4 159 & 2 000 \\
\end{longtable}}

Most targets are reached by every competitor: at five nodes and 2σ, the grammar covers them whatever the vocabulary. A witness with no constant of K₇ makes this explicit. With an empty register, only the universal leaves and operations of the grammar remain (1 to 9, π, e, φ, γ, ln 2, the ζ values), and they reach 30 of the 33 targets, all but α⁻¹, m\textsubscript{τ}/m\textsubscript{e} and m\textsubscript{μ}/m\textsubscript{e}. Three scores must therefore be kept apart: the published formulas, the free search with K₇'s alphabet, and this witness. The targets that discriminate are the few with very small relative uncertainties. α⁻¹ and m\textsubscript{μ}/m\textsubscript{e} are reached by nobody, K₇ included. Every competitor reaches the other 30 targets that K₇ reaches except m\textsubscript{τ}/m\textsubscript{e}, which K₇ and about a third of the competitors reach; the counts of the family table are exactly the counts on that one target. Ω\textsubscript{c}/Ω\textsubscript{m} shows the difference between the scores: its published value 27/32 lies outside the window of the chains (+2.46σ, main §4), yet the target is reached by K₇ with another expression, as by every competitor and by the witness. A published value outside K₇'s five-node set (α⁻¹, m\textsubscript{μ}/m\textsubscript{e} and three others) is not available to the blind search, which may still reach the target with another expression.

\subsection{S4.2.3 Comparison with the earlier runs}\label{s4.2.3-comparison-with-the-earlier-runs}

\begin{longtable}[]{@{}
  >{\raggedright\arraybackslash}p{(\linewidth - 6\tabcolsep) * \real{0.2500}}
  >{\raggedright\arraybackslash}p{(\linewidth - 6\tabcolsep) * \real{0.2500}}
  >{\raggedright\arraybackslash}p{(\linewidth - 6\tabcolsep) * \real{0.2500}}
  >{\raggedright\arraybackslash}p{(\linewidth - 6\tabcolsep) * \real{0.2500}}@{}}
\toprule\noalign{}
\begin{minipage}[b]{\linewidth}\raggedright
Quantity
\end{minipage} & \begin{minipage}[b]{\linewidth}\raggedright
first run (2026-10-17, dataset of v3.5, one-sided windows)
\end{minipage} & \begin{minipage}[b]{\linewidth}\raggedright
revision 1 (2026-10-24, refreshed dataset, two-sided windows; historical)
\end{minipage} & \begin{minipage}[b]{\linewidth}\raggedright
revision 2 (2026-11-03, current dataset)
\end{minipage} \\
\midrule\noalign{}
\endhead
\bottomrule\noalign{}
\endlastfoot
K₇ S\textsubscript{cover} & 32/33 & 31/33 & 31/33 \\
K₇ S\textsubscript{cost} (bits) & 395.6 & 409.2 & 408.1 \\
F-BETTI: reach as many as K₇ & 4 159 of 4 159 (P = 1.000) & 1 085 of 4 159 (P = 0.261) & 1 085 of 4 159 (P = 0.261) \\
F-VOCAB: reach as many as K₇ & 1 635 of 2 000 (P = 0.818) & 1 015 of 2 000 (P = 0.508) & 1 015 of 2 000 (P = 0.508) \\
F-BETTI-FICHE: reach as many as K₇ & 4 160 of 4 160 (P = 1.000) & 723 of 4 160 (P = 0.174) & 723 of 4 160 (P = 0.174) \\
P\textsubscript{global}, verdict & 1.000, LEE-RETRAIT & 0.508, LEE-RETRAIT & 0.508, LEE-RETRAIT \\
targets reached by nobody & A04 & A04, A07 & A04, A07 \\
\end{longtable}

The first run, on the dataset of version 3.5 with one-sided windows, had every alternative Betti pair at K₇'s score. Two targets change. m\textsubscript{μ}/m\textsubscript{e}, whose uncertainty in the 3.5 dataset (0.02) was some four thousand times the CODATA value (4.6 × 10⁻⁶), is now reached by nobody, K₇ included, which lowers K₇'s score by one and every competitor's alike. m\textsubscript{τ}/m\textsubscript{e}, reached by every alternative Betti pair and most invented vocabularies in the first run, is now reached by K₇ and by about a third of the competitors, its uncertainty having shrunk from 0.23 to 0.18. It is the only target that separates competitors: those that reach as many targets as K₇ are exactly those that reach m\textsubscript{τ}/m\textsubscript{e} (S4.2.2). Revision 2 changes the windows of six ratios that are not discriminating: the counts of competitors reaching as many as K₇ are therefore those of revision 1, while the costs move (K₇, 409.2 to 408.1 bits). The figures of revision 1 are historical. The verdict does not change.

\subsection{S4.2.4 What the test does not cover}\label{s4.2.4-what-the-test-does-not-cover}

The counts are descriptive. They record an absence of discrimination in these families and windows, not a measurement of statistical power. P\textsubscript{global} = 0.51 is not a probability that K₇ is right or wrong: with 30 targets reached by everyone, two by no one and one that separates, the score has little discriminating power in this protocol, which is enough to withdraw the earlier statistical claims and no more. The counts are not a physical refutation of K₇. The test does not vary the gauge group or the holonomy beyond the Betti numbers (there is no public rule for re-deriving the register in that case), does not consider other forms of an observable (an angle against its sine squared) except for the unit of δ\textsubscript{CP}, and does not look beyond five nodes, where several of K₇'s formulas lie. These freedoms remain in the residue register (main §3.3: R4, R5, R6). Nor does the test use the literal register of main §3.1 (S4.2.1).

\section{S4.3 Descriptive diagnostics of version 3.5 (archived)}\label{s4.3-descriptive-diagnostics-of-version-3.5-archived}

Version 3.5 reported the following analyses of its formula set as evidence against cherry-picking. They were computed on the formula set and dataset of version 3.5 and have not been recomputed. They describe how the formulas use the constants; they do not price the choice of formulas, which is what S4.2 does, and they are no longer read as evidence.

\begin{itemize}
\tightlist
\item
  \textbf{Complexity against deviation.} A random-forest regression of \textbar deviation\textbar{} on formula features (number of constants, largest constant, number of operations, sector), with leave-one-out cross-validation, gave R² = −0.518: formula complexity did not predict the size of the deviation.
\item
  \textbf{Sensitivity.} Perturbing each of the 20 constants by ±1 gave a 20 × 33 sensitivity matrix with 19 non-zero singular values (3.50, 2.54, 2.39, 2.14, 2.11, \ldots) and a correlation of −0.083 between sensitivity and deviation.
\item
  \textbf{Effective rank.} The 20 × 33 binary constant-usage matrix had an effective rank of 15.53; 12 singular values carried 90 \% of the variance and 17 carried 99 \%. Version 3.5 divided 33 by 15.53 and called the ratio, 2.13, an ``overdetermination ratio'' characteristic of a system that is not fitted. The ratio counts the 33 relations as independent, which they are not (main §4), and a usage matrix cannot tell a fitted formula from an unfitted one; the reading is withdrawn.
\item
  \textbf{Shared constants.} Of the 528 pairs of observables, 155 share at least one structural constant, and the 33 form one connected component; 51 pairs of deviations have \textbar ρ\textbar{} ≥ 0.5.
\end{itemize}

\section{S4.4 Coincidence probabilities of version 3.4 (archived)}\label{s4.4-coincidence-probabilities-of-version-3.4-archived}

Version 3.4 reported joint coincidence probabilities of 10⁻³⁴⁶ (uniform null) and 10⁻¹³³ (algebraic null) for the 33 Type I relations. Both are withdrawn. The uniform null treats the predictions as random numbers in {[}0, 50 \%{]} of deviation, which no formula search produces. The algebraic null compares K₇'s \emph{selected} formulas with \emph{random} formulas built from the same constants, so it does not price the selection: a competitor allowed to select its own formulas, as in S4.2, does as well. The tables are kept as a record of what was reported.

\begin{longtable}[]{@{}llll@{}}
\toprule\noalign{}
Uniform null, deviations in {[}0, 50 \%{]} & Statistic & df & p-value \\
\midrule\noalign{}
\endhead
\bottomrule\noalign{}
\endlastfoot
χ² & 1063 & 33 & 5.0 × 10⁻²⁰² \\
Fisher combined & 1100.5 & 66 & 2.2 × 10⁻¹⁸⁷ \\
KS (pull normality) & 0.189 & n/a & 0.165 \\
Combined & n/a & n/a & 10⁻³⁴⁶ \\
\end{longtable}

\begin{longtable}[]{@{}
  >{\raggedright\arraybackslash}p{(\linewidth - 4\tabcolsep) * \real{0.3333}}
  >{\raggedright\arraybackslash}p{(\linewidth - 4\tabcolsep) * \real{0.3333}}
  >{\raggedright\arraybackslash}p{(\linewidth - 4\tabcolsep) * \real{0.3333}}@{}}
\toprule\noalign{}
\begin{minipage}[b]{\linewidth}\raggedright
Algebraic null (4 188 086 formula values from the 20 constants)
\end{minipage} & \begin{minipage}[b]{\linewidth}\raggedright
K₇ (3.4 dataset)
\end{minipage} & \begin{minipage}[b]{\linewidth}\raggedright
Random formulas (median)
\end{minipage} \\
\midrule\noalign{}
\endhead
\bottomrule\noalign{}
\endlastfoot
Mean deviation & 0.73 \% & 4.1 × 10⁹ \% \\
Matches within 0.01 \% & 5 & 0 \\
Matches within 1 \% & 28 & 0 \\
\end{longtable}

The pull distribution of version 3.4 had mean −0.774 and standard deviation 5.62 (reduced χ² = 32.2, driven by δ\textsubscript{CP} and α\textsubscript{s}); these figures refer to the 3.4 dataset and are superseded by the comparisons of main §4.

\section{S4.5 Reproducibility}\label{s4.5-reproducibility}

\begin{itemize}
\tightlist
\item
  \textbf{Scorecard (S4.1).} {[}M, §4.3{]}, with the frozen inputs {[}S-frozen{]}.
\item
  \textbf{Look-elsewhere null (S4.2).} Producer \texttt{canonical/\allowbreak{}scripts/\allowbreak{}k7\_\allowbreak{}v3\_\allowbreak{}6\_\allowbreak{}lee\_\allowbreak{}null.\allowbreak{}py} (Sieve engine at commit \texttt{601283f}, grammar sha-256 pinned by the producer, depth 5, τ\textsubscript{z} = 2, seed 20261017); result \texttt{canonical/\allowbreak{}results/\allowbreak{}k7\_\allowbreak{}v3\_\allowbreak{}6\_\allowbreak{}lee\_\allowbreak{}null.\allowbreak{}json} (10 319 vocabularies, 2 h 48 on three processes); ties of the dual cost, \texttt{canonical/\allowbreak{}scripts/\allowbreak{}k7\_\allowbreak{}v3\_\allowbreak{}6\_\allowbreak{}lee\_\allowbreak{}null\_\allowbreak{}closeout.\allowbreak{}py}. The tables of S4.2 and the figures of main §7.2 are generated and checked by \texttt{canonical/\allowbreak{}scripts/\allowbreak{}k7\_\allowbreak{}v3\_\allowbreak{}6\_\allowbreak{}s4\_\allowbreak{}tables.\allowbreak{}py}. The first run and revision 1 are the same producer before each amendment, each at its own commit, read from git by the generator. The producer, the closeout, the dataset as read by the run, the preregistration with its two amendments, the result and the results of the two earlier runs are deposited with their hashes in the public repository (Arithmon/K7, \texttt{publications/\allowbreak{}validation/\allowbreak{}v3.\allowbreak{}6/\allowbreak{}}, which keeps the \texttt{canonical/\allowbreak{}} paths); the producer runs from there against a clone of the Sieve.
\item
  \textbf{Archives (S4.3, S4.4).} The scripts of the 3.4 and 3.5 releases; they are not maintained.
\end{itemize}

\section{Author's Related Works}\label{authors-related-works}

\begin{itemize}
\tightlist
\item
  \textbf{{[}M{]}} B. de La Fournière, ``A reproducible standard for arithmetic coincidence claims in fundamental physics,'' Zenodo \href{https://doi.org/10.5281/zenodo.20728964}{10.5281/zenodo.20728964} (2026). The Sieve methodology and its four-null scorecard for the K₇ relations.
\item
  \textbf{{[}S-frozen{]}} B. de La Fournière, ``Counting Coincidences: frozen inputs'' (observable list and expression grammars), Zenodo dataset \href{https://doi.org/10.5281/zenodo.20666878}{10.5281/zenodo.20666878}, v1.0 of 12 June 2026.
\end{itemize}

\begin{center}\rule{\linewidth}{0.5pt}\end{center}

\emph{The K₇ Framework: Supplement S4: Sieve and Look-Elsewhere Diagnostics}


\bigskip
% GÉNÉRÉ par gen_bib_3_6.py depuis physics_main.md : ne pas éditer.

\begin{thebibliography}{99}

\bibitem{1} Particle Data Group, F. Takahashi et al., Int. J. Mod. Phys. A 41, 2630011 (2026)

\bibitem{2} S. Weinberg, Phys. Rev.~D 13, 974 (1976)

\bibitem{3} Planck Collaboration, A\&A 641, A6 (2020); chains: Planck Legacy Archive, COM\_CosmoParams\_base-plikHM-TTTEEE-lowl-lowE\_R3.00

\bibitem{4} A.G. Riess et al., ApJL 934, L7 (2022)

\bibitem{5} C.D. Froggatt, H.B. Nielsen, Nucl. Phys. B 147, 277 (1979)

\bibitem{6} Y. Koide, Lett. Nuovo Cim. 34, 201 (1982)

\bibitem{7} C. Furey, \emph{Standard Model Physics from an Algebra?}, PhD thesis, U. Waterloo (2015)

\bibitem{8} C. Furey, M.J. Hughes, Phys. Lett. B 831, 137186 (2022)

\bibitem{9} R. Wilson, arXiv:2404.18938 (2024)

\bibitem{10} T.P. Singh et al., arXiv:2206.06911v3 (2024)

\bibitem{11} B.S. Acharya, S. Gukov, Phys. Rep.~392, 121 (2004)

\bibitem{12} L. Foscolo et al., Duke Math. J. 170, 3 (2021)

\bibitem{13} A. Kovalev, ``Coassociative K3 fibrations of compact G₂-manifolds,'' arXiv:math/0511150 (2005)

\bibitem{14} M. Haskins, J. Nordstrom, arXiv:1809.09083 (2022)

\bibitem{15} G. Bera, ``Associative submanifolds in twisted connected sum G₂-manifolds,'' arXiv:2209.00156 (2022)

\bibitem{16} D. Crowley, S. Goette, J. Nordstrom, Invent. Math. (2025)

\bibitem{17} T. Dray, C.A. Manogue, \emph{The Geometry of the Octonions}, World Scientific (2015)

\bibitem{18} J.F. Adams, \emph{Lectures on Exceptional Lie Groups}, U. Chicago Press (1996)

\bibitem{19} D.J. Gross et al., Nucl. Phys. B 256, 253 (1985)

\bibitem{20} D.D. Joyce, \emph{Compact Manifolds with Special Holonomy}, Oxford U. Press (2000)

\bibitem{21} A. Kovalev, J. Reine Angew. Math. 565, 125 (2003)

\bibitem{22} A. Corti et al., Duke Math. J. 164, 1971 (2015)

\bibitem{23} R. Harvey, H.B. Lawson, Acta Math. 148, 47 (1982)

\bibitem{24} B.S. Acharya, Class. Quant. Grav. 19, 5619 (2002)

\bibitem{25} B.S. Acharya, E. Witten, arXiv:hep-th/0109152 (2001)

\bibitem{26} The NOvA Collaboration, The T2K Collaboration, ``Joint neutrino oscillation analysis from the T2K and NOvA experiments,'' Nature 646, 818-824 (2025), doi:10.1038/s41586-025-09599-3, arXiv:2510.19888

\bibitem{27} I. Esteban, M.C. Gonzalez-Garcia, M. Maltoni, I. Martinez-Soler, J.P. Pinheiro, T. Schwetz, ``NuFit-6.0: updated global analysis of three-flavor neutrino oscillations,'' JHEP 12 (2024) 216, arXiv:2410.05380; NuFIT 6.1: three-neutrino fit based on data available in November 2025, www.nu-fit.org

\bibitem{28} L. de Moura, S. Ullrich, CADE 28, 625 (2021)

\bibitem{29} mathlib Community, github.com/leanprover-community/mathlib4

\bibitem{30} DUNE Collaboration, FERMILAB-TM-2696 (2020)

\bibitem{31} DUNE Collaboration, arXiv:2103.04797 (2021)

\bibitem{32} E. Witten, ``Strong coupling expansion of Calabi-Yau compactification,'' Nucl. Phys. B 471, 135 (1996)

\bibitem{33} B.S. Acharya et al., Phys. Rev.~D 76, 126010 (2007)

\bibitem{34} M. Atiyah, E. Witten, Adv. Theor. Math. Phys. 6, 1 (2002)

\bibitem{35} G. Kane, \emph{String Theory and the Real World} (2017)

\bibitem{36} J. Distler, S. Garibaldi, Commun. Math. Phys. 298, 419 (2010)

\bibitem{37} J.C. Baez, ``Octonions and the Standard Model,'' math.ucr.edu/home/baez/standard/ (2020-2025)

\bibitem{38} Springer Nature, ``Artificial intelligence (AI) policy,'' www.springernature.com/gp/policies (2024)

\bibitem{39} J.A. Wheeler, ``Information, physics, quantum: the search for links,'' in \emph{Complexity, Entropy, and the Physics of Information} (W.H. Zurek, ed.), Addison-Wesley (1990), pp.~3--28.

\bibitem{40} J. Worrall, ``Structural realism: the best of both worlds?,'' \emph{Dialectica} \textbf{43}, 99--124 (1989).

\bibitem{41} J. Ladyman, D. Ross, \emph{Every Thing Must Go: Metaphysics Naturalized}, Oxford University Press (2007).

\bibitem{42} R. Pinčák, A. Pigazzini, M. Pudlák, E. Bartoš, ``Geometric origin of a stable black hole remnant from torsion in G₂-manifold geometry,'' Gen.~Rel. Grav. \textbf{58}, 29 (2026), doi:10.1007/s10714-026-03528-z

\bibitem{43} D. Joyce, S. Karigiannis, ``A new construction of compact torsion-free G₂-manifolds by gluing families of Eguchi-Hanson spaces,'' \emph{J. Differential Geom.} (2021), arXiv:1707.09325 (2017).

\bibitem{44} S. Mukai, ``Finite groups of automorphisms of K3 surfaces and the Mathieu group,'' \emph{Invent. Math.} \textbf{94}, 183--221 (1988).

\bibitem{45} A. Garbagnati, A. Sarti, ``Symplectic automorphisms of prime order on K3 surfaces,'' \emph{J. Algebra} \textbf{318}, 323--350 (2007), arXiv:math/0603742.

\bibitem{46} LEP2 SUSY Working Group (ALEPH, DELPHI, L3, OPAL), ``Combined LEP chargino results, up to 208 GeV for large m0,'' note LEPSUSYWG/01-03.1 (2001), lepsusy.web.cern.ch.

\bibitem{47} Fermilab News, ``Fermilab marks major milestone for world-leading DUNE experiment,'' 7 May 2026, news.fnal.gov.

\bibitem{48} J. Chen, H. Hong, ``Intermediate curvature and splitting theorem,'' arXiv:2604.26529 (2026). (Cited in Supplement S1, §8.4.)

\bibitem{49} M. Tristram et al., ``Cosmological parameters derived from the final Planck data release (PR4),'' A\&A 682, A37 (2024). (Cited in Supplement S3; not used in the dataset.)

\bibitem{50} E. Tiesinga, P.J. Mohr, D.B. Newell, B.N. Taylor, ``CODATA recommended values of the fundamental physical constants: 2022,'' Rev.~Mod. Phys. 97, 025002 (2025). (Cited in Supplement S3.)

\bibitem{51} V. V. Nikulin, ``Integer symmetric bilinear forms and some of their geometric applications,'' Izv. Akad. Nauk SSSR Ser. Mat. 43 (1979); English transl. Math. USSR Izv. 14 (1980).

\bibitem{52} Belle II Collaboration, ``Measurement of the τ-lepton mass with the Belle II experiment,'' Phys. Rev.~D 108, 032006 (2023), doi:10.1103/PhysRevD.108.032006, arXiv:2305.19116.

\bibitem{53} S. K. Donaldson, ``Adiabatic limits of co-associative Kovalev--Lefschetz fibrations,'' in Algebra, Geometry, and Physics in the 21st Century, Progr. Math. 324, Birkhäuser (2017), pp.~1--29, arXiv:1603.08391.

\bibitem{54} C. Camere, ``Lattice polarized irreducible holomorphic symplectic manifolds,'' Ann. Inst. Fourier 66 (2016), 687--709, arXiv:1312.3224.

\bibitem{55} C. Zhou, Z. Zhou, ``Algebraic Stability and Cosmological Structure E: The Arithmetic Necessity of Three Generations of Fermions,'' Zenodo, doi:10.5281/zenodo.18626667 (2026).

\bibitem{56} Particle Data Group, S. Navas et al., ``Review of Particle Physics,'' Phys. Rev.~D 110, 030001 (2024), doi:10.1103/PhysRevD.110.030001 (cited in S3 only)

\bibitem{57} CDF Collaboration, T. Aaltonen et al., ``High-precision measurement of the W boson mass with the CDF II detector,'' Science 376, 170 (2022), doi:10.1126/science.abk1781

\bibitem{58} JUNO Collaboration, ``First measurement of reactor neutrino oscillations at JUNO,'' arXiv:2511.14593 (2025).

\bibitem{59} JUNO Collaboration, ``Sub-percent precision measurement of neutrino oscillation parameters with JUNO,'' Chin. Phys. C 46, 123001 (2022), arXiv:2204.13249.

\bibitem{60} D. Buttazzo, G. Degrassi, P.P. Giardino, G.F. Giudice, F. Sala, A. Salvio, A. Strumia, ``Investigating the near-criticality of the Higgs boson,'' JHEP 12 (2013) 089, arXiv:1307.3536.

\end{thebibliography}


\end{document}
